\documentclass[12pt,a4paper]{article}
\usepackage[margin=2.3cm]{geometry}
\usepackage{fontspec}
\setmainfont{Times New Roman}
\usepackage{setspace}
\singlespacing
\usepackage{enumitem}
\setlist{nosep,leftmargin=1.2em}
\usepackage{titlesec}
\titleformat{\section}{\normalfont\large\bfseries}{}{0em}{}
\titleformat{\subsection}{\normalfont\normalsize\bfseries}{}{0em}{}
\titlespacing{\section}{0pt}{10pt}{4pt}
\titlespacing{\subsection}{0pt}{7pt}{2pt}
\usepackage{booktabs}
\usepackage{array}
\usepackage{hyperref}
\hypersetup{hidelinks}
\setlength{\parindent}{0pt}
\setlength{\parskip}{4pt}
\newenvironment{refs}{\begin{list}{}{\setlength{\leftmargin}{1.5em}\setlength{\itemindent}{-1.5em}\setlength{\itemsep}{1pt}}}{\end{list}}

\begin{document}

{\large\bfseries BE451 Thesis Synopsis --- Working draft with seven alternative topics}

7 September 2026. Each topic below follows the course's synopsis structure (research question, closest literature and contribution, data and feasibility, empirical method, three-month scope, references) and is written to fit 1--2 pages. Pick two, rewrite them in your own words, fill the [insert] fields, and delete the rest. All references are in approved journals and verified against the publisher record.

\subsection*{Overview and ranking}

\begin{tabular}{@{}p{0.5cm}p{7.2cm}cccc@{}}
\toprule
 & Topic & Rubric/20 & Ident. & Feas. & Prize \\
\midrule
I & Corporate tax cut and interest cap 2019 & 19 & 8 & 8 & 7--8 \\
F & Investment screening laws and SME ownership changes & 19 & 6 & 7 & 8 \\
H & Profit-cap threat and welfare firms 2016--2026 & 18 & 7 & 8 & 8 \\
J & Employment protection reform 2022 and capital structure & 18 & 8 & 9 & 7 \\
G & PE roll-ups below merger thresholds and local market power & 17 & 6 & 6 & 9 \\
B & Electricity price zones 2022 and corporate liquidity & 17 & 5 & 6 & 6 \\
A & Acquisitions of Swedish SMEs by parent type & 15 & 5 & 6 & 7 \\
\bottomrule
\end{tabular}

Rubric: fit to the pattern of the 25 past prize winners (anchor paper, Swedish institutional edge, mechanism test, informative null, topical hook, own data). Ident./Feas./Prize: adversarial reviewers' scores 1--10. Recommended pair: H + I; J first reserve; G highest news value but highest data risk; F best international option. Decision rule at the first Serrano login: if interest expense, tax expense or group parent are missing, I falls and J replaces it; if municipality is missing, G and B fall.

\newpage

\section{Topic H: What Does the Threat of a Profit Cap Cost? Political Uncertainty and Sweden's Private Welfare Firms, 2016--2026}

\subsection{Tentative research question}

\emph{How did Sweden's private welfare firms respond to the threat of a statutory profit cap---in investment, payouts, entry and exit, and ownership changes---and did the response reverse when the threat was removed?}

Between November 2016 and June 2018 every private provider of tax-financed schooling, social care and personal assistance in Sweden operated under a credible legislative threat. The Reepalu inquiry (SOU 2016:78, 8 November 2016) proposed to cap operating profit at the government borrowing rate plus seven percentage points times operating capital; the government's bill (prop.\ 2017/18:159, 19 March 2018) adopted the cap; the Riksdag rejected it on 7 June 2018. The threat has returned: a directive on a school sector ``without profit motive'' (dir.\ 2022:102, 30 June 2022), an inquiry (SOU 2025:37), a bill restricting value transfers in the school sector (prop.\ 2025/26:292, 25 June 2026), and the general election of 13 September 2026 in which profit caps are a campaign issue. Private providers deliver a substantial share of Swedish schooling and care [insert share], many owned by private equity. To our knowledge no study examines the episode at the firm level.

The mechanism is the real-option value of waiting under political uncertainty (Julio and Yook, 2012; Jens, 2017): firms postpone irreversible investment until policy is resolved. A profit cap adds a sharper prediction: capital retained inside a capped firm is worth less, so owners should accelerate dividends and group contributions while the cap is pending and reverse after 7 June 2018. The competing hypothesis is that owners assigned the proposal a low probability, since it lacked a parliamentary majority, so nothing happened. A null is therefore informative: it measures how much a politically salient but legislatively weak threat actually costs.

\subsection{Closest literature and contribution}

\textbf{Julio and Yook (2012, JF)} document, across 48 countries, that firms reduce investment by 4.8 percent on average in national election years and recover afterwards. We replicate the event-window design on a single, well-dated policy threat and extend it by identifying the treated group within one country and adding payout and ownership outcomes. \textbf{Jens (2017, JFE)} shows the effect is causal using staggered US gubernatorial elections with placebo tests; we replicate the placebo logic with health-care providers (SNI 86), which the bill excluded, as a within-sector control. \textbf{Kelly, P\'astor, and Veronesi (2016, JF)} show political uncertainty is priced in option markets; we use their theory for the payout prediction and for an illustrative event study on the four listed welfare firms (Attendo, Humana, AcadeMedia, Ambea). \textbf{Bonaime, Gulen, and Ion (2018, JFE)} show policy uncertainty reduces merger activity; we use this for ownership changes.

Our contribution is threefold. First, the first firm-level evidence on how a sector-specific profit-cap threat affects investment, using the full population of Swedish private welfare firms [insert N]. Second, a test of the real-option channel against the payout channel, which make opposite predictions for dividends, with private-equity ownership as a mechanism split. Third, hand-coded legislative dates, sector scope and ownership of the largest welfare groups, which allow a placebo that election-based designs cannot provide. Benchmark: with several thousand treated firms the minimum detectable effect on investment/assets is well below the 4.8 percent of Julio and Yook, so a null rejects that magnitude in our setting.

\subsection{Data and feasibility}

Serrano (Swedish House of Finance), annual accounts of all Swedish limited companies 2012--2024: investment, dividends, group contributions, cash, employees, entry, exit, group parent; SWAMID access, dividend and appropriation fields to be confirmed at first login. Treated: private firms in SNI 85 (education) and 87--88 (care, personal assistance); placebo: SNI 86; controls: matched private service firms. Listed firms from LSEG/FinBas. Expected sample: several thousand treated firms, tens of thousands of controls. Risks and fallback: owner and management screening rules from 1 January 2019 (main window 2016--2018; SNI 85 separated from 87--88); the 2022--2026 episode coincides with the rate shock (matching, SNI 86 placebo); the listed sample is small and Ambea bid for Humana on 29 June 2026 (event study illustrative only).

\subsection{Empirical method and three-month scope}

Event-study difference-in-differences with firm and industry-by-year fixed effects around 8 November 2016, 19 March 2018 and 7 June 2018; outcomes investment/assets, dividends/equity, group contributions, employment growth, exit and ownership change. Identification: the treated set is defined by the bill's own scope; SNI 86 shares public financing and demand structure but was excluded; matched service firms absorb aggregate shocks. Mechanism: private-equity-owned versus other treated firms; firms with profits above the proposed cap versus below. Robustness: SNI 86, pre-2016 pseudo-events, exclusion of asylum accommodation (SNI 87.9), alternative matching. MVP: 2014--2020 difference-in-differences on investment and payouts with the SNI 86 placebo; extensions: the 2022--2026 school episode, the PE mechanism, the listed event study.

\subsection*{References}
\begin{refs}
\item Bonaime, Alice, Huseyin Gulen, and Mihai Ion, 2018, Does policy uncertainty affect mergers and acquisitions?, \emph{Journal of Financial Economics} 129, 531--558.
\item Jens, Candace E., 2017, Political uncertainty and investment: Causal evidence from U.S. gubernatorial elections, \emph{Journal of Financial Economics} 124, 563--579.
\item Julio, Brandon, and Youngsuk Yook, 2012, Political uncertainty and corporate investment cycles, \emph{Journal of Finance} 67, 45--83.
\item Kelly, Bryan, \v{L}ubo\v{s} P\'astor, and Pietro Veronesi, 2016, The price of political uncertainty: Theory and evidence from the option market, \emph{Journal of Finance} 71, 2417--2480.
\end{refs}

\newpage

\section{Topic I: Do Taxes Depress or Stimulate Borrowing? Sweden's 2019 Corporate Tax Cut and Interest Deduction Limit}

\subsection{Tentative research question}

\emph{When Sweden cut the corporate tax rate and capped interest deductibility on 1 January 2019, did private firms borrow less, as the trade-off theory of the tax shield predicts, or more, as the productivity channel of Ivanov, Pettit, and Whited (2026) predicts---and does the answer differ between firms that lost deductions and firms that only gained a lower rate?}

The reform (prop.\ 2017/18:245) contained two treatments with opposite predictions under standard theory. The corporate tax rate fell from 22 to 21.4 percent in 2019 and to 20.6 percent in 2021, lowering the value of the interest tax shield for every firm. At the same time, net interest expense became deductible only up to 30 percent of taxable EBITDA, with a simplification rule allowing SEK 5 million of net interest per group of associated companies regardless. Trade-off theory (Heider and Ljungqvist, 2015) predicts lower leverage from both. Ivanov, Pettit, and Whited (2026) find the opposite for small US private firms: after tax cuts, higher after-tax cash flows move default further away, credit spreads fall, and leverage rises. Sanati and Beyhaghi (2025) use the USD 25 million exemption in the 2017 US reform in a regression discontinuity and find that losing the deduction lowers debt and investment. Sweden lets us run both tests in one population---all limited companies---with the SEK 5 million threshold as a sharp discontinuity and the rate cut as a common shock whose bite varies with pre-reform taxable profits. A null on both margins is informative: Swedish private-firm capital structure would then not be tax-driven at the elasticities estimated in the US literature.

\subsection{Closest literature and contribution}

\textbf{Ivanov, Pettit, and Whited (2025, REStud)} use US state tax changes and private-firm data and find that leverage rises after tax cuts because lower taxes raise capital productivity and lower credit spreads; we replicate with the 2019 and 2021 rate steps and add the simultaneous deduction cap. \textbf{Sanati and Beyhaghi (2025, RFS)} exploit the USD 25 million exemption of the 2017 US reform in a regression discontinuity; we replicate the design at Sweden's SEK 5 million threshold measured per group. \textbf{Heider and Ljungqvist (2015, JFE)} estimate that a one percentage point tax increase raises leverage by about 0.4 percentage points; we use this as the benchmark. \textbf{Brok (2026, JFQA; 2024, Review of Finance)} shows that tax avoidance crowds out the debt tax shield (2026) and that group structure adds leverage incentives via deductible subsidiary losses (2024); we use Serrano's group identifiers to measure the threshold per group and test reallocation of interest across entities.

Our contribution is threefold: the first evidence on Sweden's 2019 reform for the full private-firm population [insert N]; a horse race between two channels with opposite signs in the same data; and a hand-built group-level net-interest variable that assigns the SEK 5 million threshold correctly, which firm-level data cannot. Benchmark: the Heider--Ljungqvist elasticity implies a leverage change of roughly 0.6 percentage points for the 1.4 point rate cut; our minimum detectable effect will be well below that.

\subsection{Data and feasibility}

Serrano, all Swedish limited companies 2013--2024: debt, interest expense, tax expense, EBITDA, investment, dividends, group parent and intra-group balances; interest, tax and group fields to be confirmed at first login. Treatments: (i) rate cut: pre-reform effective tax payments/assets, loss-making firms untreated; (ii) cap: net interest 2017--2018 above versus below SEK 5 million, aggregated within groups; (iii) rental-housing firms, which received a 12 percent primary deduction. Sample: hundreds of thousands of firms; the cap binds for a few thousand, mainly real estate and PE-owned. Risks and fallback: concurrent 2019 hybrid and CFC rules affect large multinationals (exclude); rates near zero 2019--2021 and rising in 2022 (industry-by-year effects, dose design); anticipation after the bill of 3 May 2018. Fallback: rate-cut analysis only.

\subsection{Empirical method and three-month scope}

Rate cut: difference-in-differences of leverage, new borrowing, interest expense/debt and investment on pre-reform tax exposure interacted with post-2019 and post-2021 indicators, with firm and industry-by-year fixed effects. Cap: regression discontinuity in group net interest at SEK 5 million with bunching analysis, plus difference-in-differences for firms bound by the 30 percent EBITDA rule versus unbound. Identification: the law fixes the rate path and the threshold; within industry and year, firms with different pre-reform exposure faced the same macro conditions. Threats: manipulation at the threshold (bunching is a finding; assign treatment on pre-reform net interest); intra-group reallocation (group aggregation); the rate cycle (industry-by-year effects; placebo 2016--2018). MVP: rate-cut difference-in-differences with a tax-exposure dose; extensions: threshold discontinuity, bunching, group channel, rental housing.

\subsection*{References}
\begin{refs}
\item Brok, Peter, 2024, Debt and taxes: The role of corporate group structure, \emph{Review of Finance} 28, 2051--2082.
\item Brok, Peter, 2026, Debt and taxes: The role of tax avoidance, \emph{Journal of Financial and Quantitative Analysis} 61, 2007--2032.
\item Heider, Florian, and Alexander Ljungqvist, 2015, As certain as debt and taxes: Estimating the tax sensitivity of leverage from state tax changes, \emph{Journal of Financial Economics} 118, 684--712.
\item Ivanov, Ivan T., Luke Pettit, and Toni M. Whited, 2026, Taxes depress corporate borrowing: Evidence from private firms, \emph{Review of Economic Studies} 93, 2715--2749.
\item Sanati, Ali, and Mehdi Beyhaghi, 2025, How does removing the tax benefits of debt affect firms? Evidence from the 2017 U.S. tax reform, \emph{Review of Financial Studies}, advance access, doi:10.1093/rfs/hhaf041.
\end{refs}

\newpage

\section{Topic F: Investment Screening Laws and the Market for Small Firms: Staggered Evidence from Northern Europe}

\subsection{Tentative research question}

\emph{Does investment screening reduce the number of ownership changes in small and medium-sized firms in protected sectors, and the share of foreign and private-equity buyers among them---and does the effect work through prohibition risk or through administrative cost?}

Sweden's screening act (2023:560) took effect on 1 December 2023 and is unusual in the EU: no size threshold, seven categories of protected activity, and a notification duty for all investors, including Swedish ones, while only non-EU buyers can be blocked. The Inspectorate of Strategic Products received 114 notifications in December 2023, 1,261 in 2024 and 1,987 in 2025, opened reviews in about 1.5 percent of cases and issued three prohibitions. The European Commission's fifth screening report singles out Sweden as the main source of EU case volume. Austria (25 July 2020, micro-firms exempt), Germany (expansion 1 May 2021), Denmark (1 July 2021) and Finland (2020) provide staggered variation; Norway and Switzerland are excluded (entry into force unclear or not yet); Iceland is a control. The literature on screening covers listed targets (Frattaroli, 2020; CFIUS studies); no paper uses the Nordic laws on deal data or studies private targets in Europe.

Mechanism: for SMEs prohibition risk is close to zero, so any effect must come through uncertainty, delay and administrative cost (Bonaime, Gulen, and Ion, 2018). Sweden's all-investors rule gives a triple difference---protected sector, post-December 2023, buyer origin (Swedish, EU, non-EU)---that separates prohibition risk from paperwork. A null means the law is a paper tiger; an effect means a chilling of ownership changes in small firms that the legislator did not intend.

\subsection{Closest literature and contribution}

\textbf{Frattaroli (2020, JFE)} studies France's 2014 decree with a difference-in-differences on listed targets and finds fewer foreign bids and entrenchment effects; we replicate the design on private SME targets across countries. \textbf{Dinc and Erel (2013, JF)} document government opposition to foreign bids in the EU-15 and its effect on completion; we use their framing of economic nationalism. \textbf{Dessaint, Golubov, and Volpin (2017, JFE)} provide the staggered cross-country template (employment protection reforms and takeovers); we apply it to screening laws. \textbf{Yilmaz (2023, JFQA)} shows foreign acquisitions affect industry peers' value; we use listed peers' returns around Swedish screening decisions as a secondary outcome.

Our contribution is threefold: the first deal-level evidence on the Nordic screening laws and on SME targets in Europe; a triple difference that separates prohibition risk from administrative cost, possible only under Sweden's all-investors rule; and a hand-built crosswalk from the seven protected-activity categories to industry codes for four countries [insert N of deals].

\subsection{Data and feasibility}

S\&P Capital IQ deal universe, target HQ in the eight countries 2015--2026 (own saved screen, 56,455 deals; 2025--2026 exported, remaining years in progress); Serrano for Swedish ownership changes and industry codes; buyer country both immediate and ultimate parent. Time to completion is not observable for private deals (87 percent have identical announcement and closing dates), so outcomes are deal counts and composition. Sector coding is the weakest link: Sweden defines protected activities functionally, not by industry code, so a narrow and a broad crosswalk are built and reported. Sanity check against ISP notification volumes. Risks: concurrent rate cycle and energy crisis (non-covered sectors in the same country, covered sectors in not-yet-treated countries); anticipation (bill spring 2023; leads); short Swedish post-period (33 months, sufficient for volumes).

\subsection{Empirical method and three-month scope}

Staggered difference-in-differences (Callaway and Sant'Anna estimator) at country by sector by year on log deal counts, foreign-buyer share and PE share, with country-by-year and sector fixed effects; event-study leads and lags; Swedish triple difference by buyer origin and target size; Austria's micro-firm exemption as a within-country discontinuity; Sweden's sector expansion of 1 October 2024 as a second event. MVP: Sweden alone with Serrano and Capital IQ; extension: the staggered design across Austria, Germany, Denmark and Finland. Fallback: Sweden only.

\subsection*{References}
\begin{refs}
\item Bonaime, Alice, Huseyin Gulen, and Mihai Ion, 2018, Does policy uncertainty affect mergers and acquisitions?, \emph{Journal of Financial Economics} 129, 531--558.
\item Dessaint, Olivier, Andrey Golubov, and Paolo Volpin, 2017, Employment protection and takeovers, \emph{Journal of Financial Economics} 125, 369--388.
\item Dinc, I. Serdar, and Isil Erel, 2013, Economic nationalism in mergers and acquisitions, \emph{Journal of Finance} 68, 2471--2514.
\item Frattaroli, Marc, 2020, Does protectionist anti-takeover legislation lead to managerial entrenchment?, \emph{Journal of Financial Economics} 136, 106--136.
\item Yilmaz, Umit, 2023, The effects of foreign acquisitions on the value of industry peers, \emph{Journal of Financial and Quantitative Analysis} 59, 1849--1885.
\end{refs}

\newpage

\section{Topic J: When Firing Gets Cheaper: Sweden's 2022 Employment Protection Reform, Leverage and Investment}

\subsection{Tentative research question}

\emph{When Sweden lowered the cost of dismissing employees in 2022, did firms increase financial leverage and investment, as the theory of operating leverage predicts, and was the effect largest for the firms whose dismissal rules changed most?}

Prop.\ 2021/22:176 was decided on 17 March 2022, entered into force on 30 June 2022 and applies from 1 October 2022. Three changes: ``saklig grund'' became ``sakliga skäl''; employment ends after the notice period even during a dispute, where previously full pay continued until judgment; and every employer may exempt three employees from the seniority order, where previously only employers with at most ten employees could exempt two. Firms with more than ten employees therefore went from zero to three exemptions, firms with at most ten from two to three. Treatment intensity jumps at ten employees, observable in Serrano for all limited companies, with accounts for 2022--2024 available. Earlier Swedish work on the 2001 two-exemption rule studies labour outcomes (sickness absence, productivity), not capital structure or investment [verify: von Below and Thoursie 2010; Bjuggren 2018].

Mechanism: Simintzi, Vig, and Volpin (2015) and Serfling (2016) show that stronger employment protection raises operating leverage and makes firms reduce financial leverage; Bai, Fairhurst, and Serfling (2020) show weaker protection raises investment and growth. Sweden 2022 is the reverse experiment: weaker protection should raise debt and investment, most for firms just above ten employees. Competing hypothesis: exemptions are rarely used because collective agreements dominate Swedish dismissals, so no real effect. A null is informative: it says the Swedish model's agreements, not the statute, set the cost of dismissal.

\subsection{Closest literature and contribution}

\textbf{Serfling (2016, JF)} uses staggered US state wrongful-discharge laws and finds that higher firing costs reduce leverage; we replicate the design with a reform that lowered firing costs and a size threshold instead of state timing. \textbf{Simintzi, Vig, and Volpin (2015, RFS)} show across countries that employment protection reduces leverage through operating leverage; we test the same substitution in one country with a sharp within-country treatment. \textbf{Bai, Fairhurst, and Serfling (2020, RFS)} show weaker protection raises investment and firm growth; we add investment and employment outcomes. \textbf{Dessaint, Golubov, and Volpin (2017, JFE)} show employment protection reduces takeover activity; we add ownership changes as a secondary outcome.

Our contribution is threefold: the first evidence on the 2022 reform's financial effects in the full population of Swedish firms [insert N]; a test of the operating-leverage substitution channel against the collective-agreement null; and a regression discontinuity at ten employees, a design the cross-country literature cannot use.

\subsection{Data and feasibility}

Serrano 2016--2024, all limited companies: employees, debt, interest expense, investment, bankruptcy. Collective agreement coverage by industry from Medlingsinstitutet (public). No external request needed; outcome window is three fiscal years. Risks: the 2022 rate shock coincides (the threshold design compares firms of nearly equal size facing the same rates); bunching at ten employees before the reform (check); some firms exempted through agreements before (attenuates); the electricity shock (industry-by-year effects, robustness by bidding zone).

\subsection{Empirical method and three-month scope}

Difference-in-differences and regression discontinuity at ten employees (size fixed in 2021) with industry-by-year fixed effects; outcomes leverage, new borrowing, investment, employment growth, staff turnover where observable, bankruptcy. Placebo: the same threshold in 2018--2021. Heterogeneity: collective agreement coverage, capital intensity. MVP: discontinuity and difference-in-differences on leverage and investment; extensions: heterogeneity by agreement coverage, the dispute-pay rule via bankruptcy data. Fallback: difference-in-differences without the threshold.

\subsection*{References}
\begin{refs}
\item Bai, John (Jianqiu), Douglas Fairhurst, and Matthew Serfling, 2020, Employment protection, investment, and firm growth, \emph{Review of Financial Studies} 33, 644--688.
\item Dessaint, Olivier, Andrey Golubov, and Paolo Volpin, 2017, Employment protection and takeovers, \emph{Journal of Financial Economics} 125, 369--388.
\item Serfling, Matthew, 2016, Firing costs and capital structure decisions, \emph{Journal of Finance} 71, 2239--2286.
\item Simintzi, Elena, Vikrant Vig, and Paolo Volpin, 2015, Labor protection and leverage, \emph{Review of Financial Studies} 28, 561--591.
\end{refs}

\newpage

\section{Topic G: Consolidation Below the Radar: Private Equity Roll-Ups and Market Power in Swedish Local Service Markets}

\subsection{Tentative research question}

\emph{When private-equity-backed platforms acquire small competitors in a local service market, do the margins of the remaining independent firms rise or fall, and do entry and exit change?}

Since 2015 PE-backed platforms have acquired hundreds of small Swedish firms in fragmented local service industries: dental and veterinary clinics, electrical and plumbing installers, accounting, physiotherapy and elderly care. In our Capital IQ extract for 2025--2026, platform add-ons are the largest acquisition channel for Swedish private targets (38 percent). Each deal is far below the merger-control thresholds (combined turnover above SEK 1 billion and at least two parties above SEK 200 million), so the resulting local concentration has never been reviewed. On 1 August 2026 a new law (2026:580) gave the Competition Authority the power to call in deals below the thresholds and to order serial acquirers to pre-notify, citing small and local markets. Whether the problem the law targets exists is an open question; we found no firm-level study in Sweden or Europe.

The market-power hypothesis (Kepler, Naiker, and Stewart, 2023; Cunningham, Ederer, and Ma, 2021) predicts that consolidation raises prices and margins for all local providers, including independents. The efficiency hypothesis (Eliason et al., 2020) predicts that platforms lower costs, squeeze independents' margins and push them out. A null implies contestable local markets and a law that addresses a problem that did not materialise.

\subsection{Closest literature and contribution}

\textbf{Kepler, Naiker, and Stewart (2023, JF)} show that US acquisitions structured to avoid antitrust review reduce competition and raise markups; we replicate the unreviewed-consolidation question where every deal is unreviewed by construction and extend from acquirers to the rivals left behind. \textbf{Wollmann (2026, JPE)} shows stealth consolidation in US dialysis raised concentration and lowered quality. \textbf{Eliason, Heebsh, McDevitt, and Roberts (2020, QJE)} show chains transfer their practices to acquired facilities; we replicate the facility-level event study and add non-acquired rivals. \textbf{Gupta, Howell, Yannelis, and Gupta (2024, RFS)} anchor the welfare-sector concern. \textbf{Aggarwal and Baxamusa (2025, JFQA)} provide recent evidence on stealth mergers.

Our contribution is threefold: the first universe of PE add-ons in Swedish local service markets from Serrano group-parent changes and Capital IQ sponsor data [insert N; 370 add-ons of Swedish private targets in 2025--2026 alone]; a test of market power against efficiency on the firms that were not acquired; and clinic-level dental prices from TLV, so that for one industry the price channel is observed directly.

\subsection{Data and feasibility}

Serrano: accounts, industry, registered municipality, group parent changes (the universe of add-ons), deregistration reason; platforms from Capital IQ sponsor status; dental prices from TLV's Tandpriskollen (export to be verified). Unit: firm-year and market-year (municipality or functional region by industry), 2012--2024; single-establishment firms in five to six industries; thousands of local markets. Municipality and deregistration fields to be confirmed at first login. Risks: registered office versus operating site (single-establishment firms); acquired units merged into platform entities (code mergers separately; main outcomes on independents); Capital IQ misses small add-ons (Serrano parent changes as universe). Fallback: dental and installation only.

\subsection{Empirical method and three-month scope}

Stacked difference-in-differences at the market level: treatment is the first platform acquisition in the market, dose the platform's share of market turnover; Callaway and Sant'Anna estimator; outcomes for remaining independents (operating margin, sales, wage cost per employee, exit split into bankruptcy and merger, entry) and for acquired units. Identification: matched control markets on pre-trends; platform entry timing driven by national roll-out (instrument: platform founding year by distance to its first acquisition). Threats: market definition (municipality, region, radius), reverse causality, selection of independents later acquired (exclude within three years). MVP: dental and installation; extensions: prices, dose-response, more industries.

\subsection*{References}
\begin{refs}
\item Aggarwal, Rajesh K., and Mufaddal Baxamusa, 2025, Stealth mergers and investment outcomes, \emph{Journal of Financial and Quantitative Analysis} 60, 2060--2087.
\item Cunningham, Colleen, Florian Ederer, and Song Ma, 2021, Killer acquisitions, \emph{Journal of Political Economy} 129, 649--702.
\item Eliason, Paul J., Benjamin Heebsh, Ryan C. McDevitt, and James W. Roberts, 2020, How acquisitions affect firm behavior and performance: Evidence from the dialysis industry, \emph{Quarterly Journal of Economics} 135, 221--267.
\item Gupta, Atul, Sabrina T. Howell, Constantine Yannelis, and Abhinav Gupta, 2024, Owner incentives and performance in healthcare: Private equity investment in nursing homes, \emph{Review of Financial Studies} 37, 1029--1077.
\item Kepler, John D., Vic Naiker, and Christopher R. Stewart, 2023, Stealth acquisitions and product market competition, \emph{Journal of Finance} 78, 2837--2900.
\item Wollmann, Thomas G., 2026, How to get away with merger: Stealth consolidation and its effects on US dialysis, \emph{Journal of Political Economy}, forthcoming, doi:10.1086/743329.
\end{refs}

\newpage

\section{Topic B: Powerless? Corporate Liquidity and the 2022 Electricity Shock in Southern Sweden}

\subsection{Tentative research question}

\emph{When electricity prices in bidding zones SE3 and SE4 rose to more than twice the level in SE1 and SE2 during 2022, did cash buffers and credit access determine which private firms cut investment and employment, and did the 2023 electricity support to firms act as constraint relief or as a cash windfall?}

Sweden has four bidding zones since 2011. Nord Pool annual averages were 63 and 66 öre per kWh in SE1 and SE2 but 137 and 162 öre in SE3 and SE4 in 2022 (2021: 43, 43, 67, 82; 2023: 46, 46, 59, 80). In 2023 the government paid support to firms only in SE3 (50 öre per kWh) and SE4 (79 öre) on consumption between October 2021 and September 2022, capped at SEK 20 million per group (SFS 2023:233). In the financing-frictions view, firms with thin cash and credit cannot smooth a transitory cost shock and cut investment and jobs, and a later cash injection restores investment in constrained firms. In the pass-through view, effects scale with electricity intensity regardless of balance sheets and a windfall raises cash and dividends (Blanchard, L\'opez-de-Silanes, and Shleifer, 1994). A null on the constraint interaction is informative: the shock was absorbed through prices, and the support was a transfer.

\subsection{Closest literature and contribution}

\textbf{Bartram, Hou, and Kim (2022, JFE)} show with California's cap-and-trade that constrained firms shift activity out of the regulated region; we replicate the regional-shock-by-constraint design with bidding zones. \textbf{Brown, Gustafson, and Ivanov (2021, JF)} show firms hit by severe weather draw on credit lines while constrained firms cut investment; we replicate the logic that responses to transitory shocks reveal frictions. \textbf{Hausman (2025, AER)} shows transmission constraints redistribute corporate profits across regions; we use the Swedish zones as the same mechanism. \textbf{Blanchard, L\'opez-de-Silanes, and Shleifer (1994, JFE)} find cash windfalls are retained or paid out; the 2023 support is our windfall test. \textbf{Almeida, Campello, and Weisbach (2004, JF)} provide the constraint measure.

Our contribution is threefold: a firm-level exposure measure from monthly zone prices and industry electricity intensity for about 100,000 firms; a test of financing frictions against pass-through with a placebo in the untreated north; and the first evaluation of the 2023 support to firms.

\subsection{Data and feasibility}

Serrano 2015--2024 (municipality of registered office and establishment count to confirm); Nord Pool monthly zone prices; SCB electricity use per unit of value added by industry; municipality-to-zone mapping. Sample: single-unit private firms with at least five employees outside groups, about 100,000 firms. Outcomes: cash, net debt, trade payables, interest coverage, dividends, shareholder contributions, investment, employment, exit. Support amounts are tax-secret at the firm level, so the windfall test uses SE4 versus SE3 intensity in 2023--2024. Fallback: Retriever addresses, or an industry-intensity design without zones.

\subsection{Empirical method and three-month scope}

Event-study difference-in-differences with firm, industry-by-year and municipality fixed effects; continuous exposure (zone price times industry intensity) interacted with year indicators 2016--2024 and with pre-determined constraint measures. Identification: within industry and year, southern firms faced higher input prices set by interconnector capacity and continental gas prices. Threats: registered office versus plant (single-unit firms); unobserved hedging (intention-to-treat); SE4 structural differences (continuous exposure, border municipalities); rate hikes interacting with leverage (constraint interaction as placebo in SE1/SE2); the gap opened mid-2021 and closed late 2022 (monthly prices). MVP: zone-by-intensity event study with constraint heterogeneity; extensions: balance-sheet transmission and the 2023 windfall chapter.

\subsection*{References}
\begin{refs}
\item Almeida, Heitor, Murillo Campello, and Michael S. Weisbach, 2004, The cash flow sensitivity of cash, \emph{Journal of Finance} 59, 1777--1804.
\item Bartram, S\"ohnke M., Kewei Hou, and Sehoon Kim, 2022, Real effects of climate policy: Financial constraints and spillovers, \emph{Journal of Financial Economics} 143, 668--696.
\item Blanchard, Olivier Jean, Florencio L\'opez-de-Silanes, and Andrei Shleifer, 1994, What do firms do with cash windfalls?, \emph{Journal of Financial Economics} 36, 337--360.
\item Brown, James R., Matthew T. Gustafson, and Ivan T. Ivanov, 2021, Weathering cash flow shocks, \emph{Journal of Finance} 76, 1731--1772.
\item Hausman, Catherine, 2025, Power flows: Transmission lines, allocative efficiency, and corporate profits, \emph{American Economic Review} 115, 2574--2615.
\end{refs}

\newpage

\section{Topic A: Buy, Decentralise, Grow? Parent Type and the Financing of Acquired Swedish Private Firms}

\subsection{Tentative research question}

\emph{When a Swedish private firm is acquired, does the type of its new parent---a private-equity-backed platform, a private equity fund, a listed decentralised serial acquirer, or an independent strategic buyer---determine whether the target's financial constraints are relieved and growth follows?}

Roughly 1,000 private Swedish firms change owner in a recorded transaction each year. In our Capital IQ extract for 2025--2026, 38 percent of these acquisitions were add-ons by PE-owned platforms, 12 percent direct fund acquisitions, 8 percent by listed strategic buyers including serial acquirers such as Lifco and Indutrade, and 42 percent by other strategic buyers. Existing Swedish evidence (Olsson and T\aa g, 2025; K\"arn\"a and Myers, 2025) compares PE-owned firms with firms never acquired, so it cannot separate the owner type from the ownership change. The mechanism is the financing synergy of Erel, Jang, and Weisbach (2015): access to a parent's internal capital market lowers cash holdings and the sensitivity of investment to cash flow and raises investment and employment. The competing hypothesis is that growth requires active integration (Eliason et al., 2020), so a decentralised parent shows no growth beyond selection. A zero difference across parent types is informative: post-deal growth would then reflect which firms get bought, not who buys them.

\subsection{Closest literature and contribution}

\textbf{Erel, Jang, and Weisbach (2015, JF)} study 5,187 acquisitions of mostly private European targets with matched before/after comparisons and find lower cash, higher investment and lower cash-flow sensitivities after the deal; we replicate on Swedish targets 2015--2023. \textbf{Boucly, Sraer, and Thesmar (2011, JFE)} find that French private-to-private LBOs raise sales by 12 percent and employment by 18 percent; we replicate and separate fund buyouts from platform add-ons. \textbf{Cohn, Hotchkiss, and Towery (2022, Review of Finance)} show PE targets are either weakly profitable or growth-constrained; we use their constraint classification to discriminate selection from treatment. \textbf{Bernstein, Lerner, and Mezzanotti (2019, RFS)} and \textbf{Olsson and T\aa g (2025, JF)} anchor the crisis and Swedish register dimensions.

Our contribution is threefold: a hand-built dataset classifying about 6,000 acquired Swedish firms into four parent types; a test of financing synergy against integration on the same population; and the first evidence on targets of decentralised serial acquirers, the model that the 2026 Competition Act amendment and the Riksbank's 2025 private equity review both single out. Benchmark: Boucly et al.\ report employment growth of 18 percent; our minimum detectable effect is around 1--1.5 percentage points per year.

\subsection{Data and feasibility}

Deals from our Capital IQ screen (56,455 DACH and Nordic transactions 2015--2026), parent type from buyer type, company type and sponsor status; outcomes from Serrano; board birth years from Retriever for a succession subsample. Deals 2015--2023, accounts 2011--2024. Expected sample: 8,000--9,000 Swedish private targets, about 6,000 matched to organisation numbers, 30,000 matched controls. Fallback: use Serrano's group-parent changes as the universe and Capital IQ only to classify the parent.

\subsection{Empirical method and three-month scope}

Event-study difference-in-differences with firm and industry-by-year fixed effects, Callaway--Sant'Anna estimator with matched controls, coefficients for years $-4$ to $+3$ by parent type; investment--cash-flow sensitivity with post and parent-type interactions. The headline is the difference between parent types, since selection on expected growth is common to all acquirers. Threats: selection on growth (pre-trends, constraint splits); accounting artefacts after group inclusion (EBIT before appropriations); attrition through merger into the parent (tracked, bounded); placebo on withdrawn deals. MVP: Erel et al.\ replication, PE versus non-PE; extension: four parent types, mechanism tables.

\subsection*{References}
\begin{refs}
\item Bernstein, Shai, Josh Lerner, and Filippo Mezzanotti, 2019, Private equity and financial fragility during the crisis, \emph{Review of Financial Studies} 32, 1309--1373.
\item Boucly, Quentin, David Sraer, and David Thesmar, 2011, Growth LBOs, \emph{Journal of Financial Economics} 102, 432--453.
\item Callaway, Brantly, and Pedro H. C. Sant'Anna, 2021, Difference-in-differences with multiple time periods, \emph{Journal of Econometrics} 225, 200--230.
\item Cohn, Jonathan B., Edith S. Hotchkiss, and Erin M. Towery, 2022, Sources of value creation in private equity buyouts of private firms, \emph{Review of Finance} 26, 257--285.
\item Eliason, Paul J., Benjamin Heebsh, Ryan C. McDevitt, and James W. Roberts, 2020, How acquisitions affect firm behavior and performance: Evidence from the dialysis industry, \emph{Quarterly Journal of Economics} 135, 221--267.
\item Erel, Isil, Yeejin Jang, and Michael S. Weisbach, 2015, Do acquisitions relieve target firms' financial constraints?, \emph{Journal of Finance} 70, 289--328.
\item Olsson, Martin, and Joacim T\aa g, 2025, What is the cost of privatization for workers?, \emph{Journal of Finance} 80, 2107--2151.
\end{refs}

\end{document}
